
\begin{UPSTIexercice}{Statut d'un capteur}
    On souhaite écrire un programme qui affiche en clair le statut d'un capteur industriel, à partir d'un code entier lu sur son bus de communication. Voici le cahier des charges :
    \begin{itemize}
        \item Le code du capteur est stocké dans une variable \texttt{int codeStatut}.
        \item Si \texttt{codeStatut} vaut \texttt{0}, afficher \texttt{Capteur a l'arret}.
        \item Si \texttt{codeStatut} vaut \texttt{1}, afficher \texttt{Capteur en fonctionnement normal}.
        \item Si \texttt{codeStatut} vaut \texttt{2} \textbf{ou} \texttt{3}, afficher \texttt{Capteur en alerte} (les deux codes correspondent à deux niveaux d'alerte différents, mais partagent le même message).
        \item Pour toute autre valeur de \texttt{codeStatut}, afficher \texttt{Code statut inconnu}.
        \item Dans tous les cas, une fois le statut affiché, le programme affiche ensuite \texttt{Fin de lecture}.
    \end{itemize}
    \UPSTIquestion{Écrire, en langage C, la structure \texttt{switch} correspondant \textbf{exactement} à ce cahier des charges. On testera le programme avec \texttt{codeStatut = 3}.}
    \UPSTIquestion{Le cahier des charges impose-t-il d'utiliser \texttt{switch} plutôt qu'une cascade de \texttt{if}/\texttt{else if} ? Justifier en une phrase à l'aide du rappel de syntaxe.}
\end{UPSTIexercice}

\begin{UPSTIprofOnlyEnv}
    \begin{UPSTIcorrectionP}{Statut d'un capteur}
        \begin{lstlisting}[language=c]
int codeStatut = 3;

switch (codeStatut)
{
    case 0:
        printf("Capteur a l'arret\n");
        break;
    case 1:
        printf("Capteur en fonctionnement normal\n");
        break;
    case 2:
    case 3:
        printf("Capteur en alerte\n");
        break;
    default:
        printf("Code statut inconnu\n");
        break;
}
printf("Fin de lecture\n");
        \end{lstlisting}
        Avec \texttt{codeStatut = 3}, l'affichage obtenu est~:
        \begin{lstlisting}[language=bash,style=console]
Capteur en alerte
Fin de lecture
        \end{lstlisting}
        Points de rigueur à vérifier dans la copie de l'étudiant~:
        \begin{itemize}
            \item les cas \texttt{2} et \texttt{3} sont bien regroupés par fall-through volontaire (\texttt{case 2:} vide, immédiatement suivi de \texttt{case 3:}), et non dupliqués ;
            \item chaque \texttt{case} se termine par un \texttt{break;}, y compris le dernier avant le \texttt{default} (bonne pratique, même si techniquement facultatif ici) ;
            \item \texttt{printf("Fin de lecture\textbackslash n");} est placé \textbf{après} l'accolade fermante du \texttt{switch}, pas à l'intérieur d'un \texttt{case} : le cahier des charges le demande \og dans tous les cas \fg, donc une seule fois, indépendamment du \texttt{case} exécuté.
        \end{itemize}

        Non, rien n'impose \texttt{switch} ici : toutes les conditions du cahier des charges sont des tests d'\textbf{égalité} (\texttt{==}, éventuellement combinés par un \og ou \fg) sur une \textbf{même variable entière}, ce qui les rend traduisibles indifféremment en \texttt{switch} ou en cascade \texttt{if}/\texttt{else if}. Le choix de \texttt{switch} est ici un choix de \textbf{lisibilité} (nombreuses valeurs testées sur la même variable), pas une obligation du langage.
    \end{UPSTIcorrectionP}
\end{UPSTIprofOnlyEnv}
